\documentclass[11pt,a4paper]{article}

% ---------------------------------------------------------
% PREAMBLE (stable, warning-free, Overleaf-safe)
% ---------------------------------------------------------
\usepackage[utf8]{inputenc}
\usepackage[T1]{fontenc}
\usepackage{lmodern}
\usepackage{amsmath, amssymb, amsfonts}
\usepackage{bm}
\usepackage{geometry}
\usepackage{graphicx}
\usepackage{hyperref}
\usepackage{cite}
\usepackage{authblk}
\usepackage{physics}
\usepackage{mathtools}

\geometry{margin=2.4cm}

\title{\textbf{Companion LXXXVII: Simulation-Based Calibration of Neural RDCC Likelihoods Using Flow-Matching and Score-Transport Diagnostics}}
\author{Michael Lehmann \\ \texttt{mi.lehmann@gmx.de}}
\date{April 2026}

\begin{document}
\maketitle

\begin{abstract}
Neural likelihoods provide a flexible and differentiable framework for combining 
topological transport statistics, Minkowski functionals, and shear power spectra 
in RDCC parameter inference. However, such likelihoods require careful 
calibration to ensure statistical correctness, unbiased parameter recovery, and 
robust uncertainty quantification. In this work, we introduce a simulation-based 
calibration pipeline for neural RDCC likelihoods using flow-matching, 
posterior-predictive checks, and score-transport diagnostics. Flow-matching 
aligns the learned likelihood with the true simulation manifold, while 
score-based diagnostics quantify discrepancies between the learned and true 
summary-statistic distributions. This Companion establishes the statistical 
foundations necessary for applying neural RDCC likelihoods to Euclid-like data.
\end{abstract}
% ---------------------------------------------------------
\section{Introduction}

Neural likelihoods offer a powerful way to model complex, non-Gaussian summary 
statistics. Companion LXXXVI introduced a unified flow-based likelihood combining 
topological transport statistics, Minkowski functionals, and the shear power 
spectrum. However, neural likelihoods must be calibrated to ensure:

\begin{itemize}
    \item unbiased parameter inference,
    \item correct posterior coverage,
    \item robustness to simulation noise,
    \item stability across the RDCC parameter space.
\end{itemize}

This Companion introduces simulation-based calibration tools tailored to RDCC.
% ---------------------------------------------------------
\section{Flow-Matching for Likelihood Correction}

Flow-matching aligns the learned likelihood $\mathcal{L}_{\theta}$ with the true 
simulation distribution by minimizing



\[
    \mathcal{L}_{\mathrm{FM}}
    =
    \mathbb{E}_{\alpha, \mathbf{s}}
    \left[
        \left\|
            f_{\theta}(\mathbf{s}) - f_{\mathrm{sim}}(\mathbf{s})
        \right\|^{2}
    \right],
\]



where $f_{\theta}$ is the learned flow and $f_{\mathrm{sim}}$ is the flow 
estimated directly from simulations.

This ensures that the learned likelihood reproduces the true simulation manifold.
% ---------------------------------------------------------
\section{Score-Transport Diagnostics}

We define the score discrepancy



\[
    \Delta s(\mathbf{s})
    =
    s_{\theta}(\mathbf{s})
    -
    s_{\mathrm{sim}}(\mathbf{s}),
\]



where $s_{\theta}$ is the score of the learned likelihood and $s_{\mathrm{sim}}$ 
is the score estimated from simulations via denoising score matching.

The norm $\|\Delta s\|$ quantifies how well the learned likelihood captures the 
geometry of the summary-statistic distribution.
% ---------------------------------------------------------
\section{Posterior-Predictive Checks}

Given observed summary statistics $\mathbf{s}_{\mathrm{obs}}$, we generate 
posterior-predictive samples



\[
    \mathbf{s}^{(i)} \sim p(\mathbf{s} \mid \mathbf{s}_{\mathrm{obs}}).
\]



We compare:

\begin{itemize}
    \item transport distances,
    \item Minkowski functional curves,
    \item shear power spectra,
\end{itemize}

between predictions and observations.

Discrepancies indicate model misspecification or insufficient training.
% ---------------------------------------------------------
\section{Coverage Tests for \texorpdfstring{$\alpha$}{alpha}}

We draw $\alpha$ from a prior, simulate summary statistics, and infer 
$\hat{\alpha}$ using the neural likelihood. Coverage is correct if



\[
    \mathbb{P}\!\left(
        \alpha \in \mathrm{CI}_{68\%}
    \right)
    \approx 0.68.
\]



This ensures that the neural likelihood yields statistically valid uncertainties.
% ---------------------------------------------------------
\section{Conclusion}

Neural likelihoods provide a powerful and flexible framework for RDCC parameter 
inference, but require careful calibration. Flow-matching aligns the learned 
likelihood with the true simulation manifold, score-transport diagnostics quantify 
geometric discrepancies, and posterior-predictive checks ensure robustness. These 
tools establish the statistical foundations necessary for applying neural RDCC 
likelihoods to Euclid-like data.

\bibliographystyle{apsrev4-2}
\nocite{*}
\bibliography{references}


\end{document}
